\scanpage{145}

\mspage{474}

pourtant plusieurs années, je crois, avant que je m'intériorise et ne reprenne à mon compte ces attitudes là, comme la honte de me montrer nu, allant de pair \unsure{ave|avec} une relation ambigüe avec mon corps. Cette honte,inculquée dès le jeune âge, est un des aspects d'une division profonde, où le corps est objet d'un tacite \unsure{mépris|mépri} alors que les valeurs dites "culturelles" (confondues avec des capacité intellectuelles de mémorisation et autres) sont montées en épingle. Cette division en m\handwritten{oi} est restée igorée jusqu'à ma quarante-huitième année, où elle a commencé à se r\handwritten{é}soudre. C'est là le deuxième grand tournant de ma vie, qui marque l'avènement de la "troisième période" dans l'histoire de ma relation à moi-même, c'est-à-dire \unsure{asi|aussi} celle de ma relation à mon corps, et à "l'homme" et à "la femme" en moi. Mais auparavant j'ai eu ample occasion de contribuer à transmettre cette division à mes enfants (*)\handwritten{,} que j'ai pu voir la transmettre à leur tour...

J'ai fait allusion déjà hier (**) au "basculement" qui a fini par avoir lieu en moi. Avec un décalage de plus de deux ans après l'arrachement au milieu familial initial (ou pour mieux dire, après la \quad\mbox{d e s t r u c t i o n}\quad de ce milieu), ce basculement consacre la mise en place des mécanismes répressifs courants, dont mon enfance avait eu la rare chance d'être exempte jusque là. J'ai détecté jusqu'à présent deux grandes forces de nature répressive, qui ont dominé ma vie d'adulte et une grande partie de mon enfance ($^{108}1$). Je crois pouvoir dir\handwritten{e} que leur apparition ne s'est pas faite progressivement, mais que dans mon cas ce\handwritten{s} mécanismes sont apparus plus ou moins du jour au lendemain et dans toute leur force, comme conséquence d'un \quad\mbox{c h o i x}\quad délibéré, au niveau inconscient. J'a\handwritten{i} qualifié précédemment ce choix d'"abdication", mais en même temps c'était aussi un puissant principe d'action: le "je serai comme "eux"" (et pas "comme moi") signifiait aussi: je vais "miser" sur "la tête", pas plus mauvaise chez moi que chez quiconque après tout, et me battre et "les" battre avec leurs propres armes

L'un de ces mécanismes, et celui qui m'interesse surtout ici, est un des plus communs qui soit: c'est la \quad\mbox{r é p r e s s i o n}\quad\mbox{d e}\quad\mbox{m e s}\quad\mbox{\unsure{t r a i t|t r a i t s}}\quad "\mbox{f é m i n i n s}" (ou ceux ressentis comme tels par les consensus courants), au profit de valeurs "viriles". L'endroit de la médaille était bien sûr l'inverstissement à fond sur mes traits et aptitudes ressentis comme "virils" et le développement à outrance de ceux-ci, qui ont pris une place démesurée.

\noindent\rule{.21\linewidth}{.35pt}

(*) Du moins, aux quatre parmi eux que j'ai contribué à élever. Le cinquième et dernier est élevé par sa mère, et jusqu'à présent il ne s'est pas présenté une occasion propice pour seulement faire connaissance, lui et moi.

(**) Voir les débuts de la note précédente "Eclosion de la force - ou les épousailles", note n° 10\deleted{\unsure{8|9}}\handwritten{\unsure{7|9}}.

% archive sheet number: 144
